\section{Audit Problem and Trust Model}\label{sec:problem}

\subsection{Completed episode and trace contract}

A completed data-agent episode is
\[
E=(q,\mathcal D,\tau,\hat a),
\]
where $q$ is the question, $\mathcal D$ is the accessible data environment, $\hat a$ is the stated answer, and $\tau$ is the final host trace. We write
\[
\tau=(P^*,F_{\mathrm{decl}},F_{\mathrm{host}},\operands,
      s_{\mathrm{host}},v_{\mathrm{host}}).
\]
Here $P^*$ is the final program snapshot, $F_{\mathrm{decl}}$ is the evidence declared by the agent, $F_{\mathrm{host}}$ contains observed data-source reads in the instrumented namespace, $\operands$ contains values reached by the final computation, and $(s_{\mathrm{host}},v_{\mathrm{host}})$ records host execution status and comparable output when available. Program versions, tool calls, exceptions, and timestamps remain in the full record but are omitted from the tuple.

A canonical resolver maps declarations into an evidence view
\[
F_{\mathrm{res}}=R(F_{\mathrm{decl}})
 =\{\operatorname{canon}(f):f\in F_{\mathrm{decl}},\ f\text{ resolves}\}.
\]
During audit, restricted replay produces
\[
\rho=(F_{\mathrm{replay}},s_{\mathrm{replay}},v_{\mathrm{replay}}),
\]
where $F_{\mathrm{replay}}$ records reads observed inside the replay sandbox. Keeping host reads and replay reads separate distinguishes an undeclared dependency in the original episode from a boundary violation during verification.

The audit output is
\[
\mathsf{Audit}(E)\in
\{\cert_{\mathrm{replay}},\cert_{\mathrm{closure}},
  \failure(c),\unknown\}.
\]
A failure records an observed violation and its witness; $\unknown$ records insufficient evidence or an unavailable runtime. This distinction prevents missing evidence from being interpreted as a detected error.

\subsection{Independent provenance and semantic statuses}\label{sec:statuses}

The provenance status is produced by \ledger from $\tau$, $R$, a restricted runtime, and a frozen comparator. Semantic validation receives a task-specific specification $V_q$ and returns
\[
S(E;V_q)\in\{\mathrm{correct},\mathrm{incorrect},\mathrm{unknown}\}.
\]
Depending on the task, $V_q$ may contain reference SQL, a normalized program specification, database variants, executable constraints, metamorphic relations, or a human judgement. The provenance and semantic statuses are stored as independent fields. Neither is used as an input to the other.

\subsection{Threat and trust assumptions}\label{sec:threat}

The host agent, generated program, stated answer, and declared citations are fallible. The trusted computing base consists of the trace logger, canonical resolver, restricted replay sandbox, and type-aware comparator. The analysis assumes that the logger captures the final program snapshot and all data-source reads in the instrumented namespace, the resolver maps identifiers consistently, and the sandbox enforces its mounted evidence view. These assumptions make the guarantee relative to recorded execution evidence.

Tampering with the logger, sandbox escape, nondeterministic external services, and dependencies outside the instrumented namespace are outside the present guarantee. Such events should yield $\unknown$ or a platform-level security failure rather than an execution-provenance certificate. This scope aligns the paper with post-execution auditing: it provides action recoverability and evidence integrity for computed answers, while lifecycle-wide accountability requires additional controls \citep{nian2026auditable}.
